\section{Testing Methodology: the Hostile Audit Suite}
\label{sec:testing}

This chapter describes how the claims of this project are verified. The
verification artifact is a single audit suite of 491 individual checks,
organized into fourteen sections named A through N, together with a separate
set of engine and server tests that cover the browser-based mission console.
The suite is not a friendly smoke test that exercises the happy path and
declares victory. It is written as a hostile auditor: its design goal is to
break every claim the project makes, and a check is only considered worth
keeping if there exists a plausible bug, regression, or exaggeration that
would cause it to fail.

\subsection{Philosophy: an auditor, not a cheerleader}

The motivating observation is that a research codebase of this kind makes
claims at several very different levels of abstraction, and each level
admits a different failure mode. At the lowest level, the QUBO construction
could be internally inconsistent: penalties could fail to dominate soft
terms, or the energy function could disagree with the $Q$ matrix that is
actually handed to the solver. At the solver level, an annealer could
report energies that look plausible while silently violating hard
constraints, or a heuristic could be credited with quality it does not
possess. At the pipeline level, the trust model, the GNN, and the
assignment rule could each be correct in isolation yet composed incorrectly.
At the reporting level, a number quoted in the results chapter could drift
away from the code that produced it as the implementation evolves. A test
suite that only checks one of these levels leaves the others unguarded.

The audit suite therefore treats every quantitative statement in this
report as a hypothesis to be attacked. Each check records three things: a
name, a pass or fail verdict, and a detail string that captures the
relevant numbers for postmortem analysis. The suite exits with a nonzero
status if any single check fails, so the cost of a regression is immediate
and binary rather than a slow decay of confidence. Two structural
conventions deserve emphasis.

First, several checks are \emph{floors} rather than point assertions. For
example, the production projected annealer is required to remain at least
90\% exact against exact ground truth on its benchmark instances, and the
MILP comparison enforces an exactness floor of 80\%. A floor does not
assert that the current value is optimal; it asserts that the current value
may not degrade. If a future refactor quietly drops the exactness rate from
95\% to 85\%, the build fails even though 85\% might still look respectable
in isolation. Floors convert historical achievements into permanent
obligations.

Second, the suite is deterministic. Every randomized check is seeded, and
the seeds are fixed constants in the suite rather than drawn from the
system clock. A failing check reproduces identically on every machine and
every run, which is essential when the failure is a 1-in-25 solver outcome
rather than a crisp logical error. Determinism also means the suite
functions as an executable specification: the exact instance distributions,
perturbation grids, and tolerance choices are readable code, not folklore.

\subsection{Ground-truth discipline}
\label{sec:groundtruth}

A solver audit is only as strong as the ground truth it compares against,
and exact ground truth for constrained combinatorial problems is
expensive. The suite resolves this with a tiered discipline in which the
verification method escalates with instance size, and every tier is
cross-checked against the tier below it before being trusted.

The foundation is an \emph{independent energy re-derivation}. The QUBO
energy is implemented twice: once through the $Q$ matrix that solvers
consume, and once as a direct from-scratch evaluation of the objective and
penalty terms on a selected edge set. Checks across the suite evaluate both
forms on identical solutions and require exact agreement. This catches the
most insidious class of bug in quadratic formulations: a $Q$ matrix that is
a valid encoding of \emph{some} energy function, just not the intended one.

On top of this foundation, four tiers of ground truth are used.
Enumeration of all $2^k$ optional-edge subsets, with a feasibility check on
each, is the gold standard and is used up to roughly 16 optional edges; at
that size it is cheap, total, and immune to solver error. For instances too
large to enumerate but with at most about 60 QUBO variables, the suite
solves an exact linearization of the QUBO with the CBC branch-and-cut
solver, and this linearized solve is itself cross-checked against
enumeration on the overlapping size range before it is trusted as truth.
Beyond 60 variables, exact truth comes from the MILP formulation of the
original constrained problem, whose single-commodity flow relaxation is
integral for integer demands and is therefore exact despite using
continuous flow variables; the MILP is verified against brute-force
enumeration on small instances for the same reason. Finally, at the largest
sizes, where no exact method is affordable, the suite falls back to
best-known solutions aggregated across multiple annealer seeds, and every
such comparison is explicitly labeled as a gap against a best-known bound
rather than against a proven optimum. We never silently mix the two
notions: a gap against a best-known solution is an upper bound on
suboptimality, and the report treats it as such.

\subsection{The fourteen sections}

Table~\ref{tab:auditsections} summarizes the fourteen sections;
Figure~\ref{fig:audittaxonomy} groups them into the five concerns they
address. We then discuss each section in turn.

\begin{table}[t]
  \centering
  \caption{The fourteen sections of the hostile audit suite.}
  \label{tab:auditsections}
  \begin{tabular}{cl}
    \toprule
    Section & Purpose \\
    \midrule
    A & Construction invariants of the QUBO encoding \\
    B & Bit-level simulated annealing against brute-force optima \\
    C & Analytic optimum checks on hand-computable instances \\
    D & Randomized scenario matrix at sizes beyond enumeration \\
    E & Validation and edge-case behavior of the constructor \\
    F & Scaling sanity across growing instance sizes \\
    G & Targeted domain scenarios drawn from mission structure \\
    H & Benchmark of external QUBO solvers on the real $Q$ matrix \\
    I & Trust-GNN layer correctness \\
    J & Optimality-gap benchmark, small and large instances \\
    K & The simulation harness itself \\
    L & Regressions from the earlier external review round \\
    M & MILP ground truth against the projected annealer \\
    N & Sensitivity diagnostics across hyperparameter perturbations \\
    \bottomrule
  \end{tabular}
\end{table}

\begin{figure}[t]
  \centering
  \begin{tikzpicture}[
      font=\small,
      root/.style={rectangle, draw, rounded corners, very thick,
        align=center, inner sep=5pt},
      group/.style={rectangle, draw, rounded corners, align=center,
        inner sep=4pt},
      leaf/.style={rectangle, draw, densely dashed, align=center,
        inner sep=3pt, font=\footnotesize},
      edge from parent/.style={draw, -{Stealth[length=2mm]}, thick},
      level 1/.style={sibling distance=26mm, level distance=38mm},
      level 2/.style={sibling distance=22mm, level distance=26mm},
    ]
    \node[root] {audit suite\\ 491 checks}
      child {node[group] {construction}
        child {node[leaf] {A}}
        child {node[leaf] {E}}}
      child {node[group] {ground truth}
        child {node[leaf] {B, C}}
        child {node[leaf] {D, J, M}}}
      child {node[group] {solvers}
        child {node[leaf] {H}}
        child {node[leaf] {N}}}
      child {node[group] {pipeline}
        child {node[leaf] {I, K}}
        child {node[leaf] {L}}}
      child {node[group] {robustness}
        child {node[leaf] {F}}
        child {node[leaf] {G}}};
  \end{tikzpicture}
  \caption{Taxonomy of the audit suite. The five branches group the
  fourteen sections by the concern they defend: correct construction of
  the encoding, verified ground truth, honest solver comparison, correct
  end-to-end pipeline behavior, and robustness of scale and domain
  assumptions.}
  \label{fig:audittaxonomy}
\end{figure}

\paragraph{A: construction invariants.}
Section~A attacks the encoding before any solver runs. It verifies that the
utility weights sum to exactly one, that every hard-constraint penalty
dominates the total achievable soft utility (so that no feasible-looking
violation can ever be profitable), that edges are canonicalized so the same
undirected link cannot appear twice under two orderings, and that no
self-loops can enter the variable set. These are the invariants every later
section silently relies on, so they are checked first and on every
instance the suite constructs.

\paragraph{B: bit-level annealing against brute force.}
Section~B pits the baseline bit-level simulated annealer against full
enumeration on small instances. This serves two purposes: it validates the
enumeration harness against a completely independent optimizer, and it
documents, with evidence, the known weakness of unconstrained bit-flip
annealing on penalized QUBOs, namely that the annealer cannot cross the
penalty barriers and returns infeasible states on most instances.

\paragraph{C: analytic optima.}
Section~C uses hand-computable instances whose optimal edge sets are known
by construction: topologies where a single spanning structure is forced,
budgets that admit exactly one affordable tree, and feature vectors that
make one edge strictly dominant. Because the optimum is derived on paper
rather than by any optimizer, these checks catch bugs that enumeration and
the solvers might share.

\paragraph{D: randomized scenario matrix.}
Section~D generates randomized instances at sizes beyond enumeration and
checks structural properties rather than exact optima: feasibility of every
returned solution, agreement between the two energy implementations, and
monotonicity expectations such as larger budgets never forcing worse soft
energy. This is the fuzzing layer of the suite.

\paragraph{E: validation and edge cases.}
Section~E enumerates the ways the constructor must fail cleanly.
Disconnected inputs, instances with zero optional edges, a single rover,
feature values outside $[0,1]$, and budgets below the cheapest spanning
tree must all raise clean, specific errors rather than produce silent
nonsense or opaque solver failures. Half of this section is a list of
expectations that certain calls must throw.

\paragraph{F: scaling sanity.}
Section~F verifies that construction time, variable counts, and solve
behavior scale as predicted as instances grow, catching accidental
quadratic or worse blowups in the encoding path before they reach an
experiment.

\paragraph{G: targeted domain scenarios.}
Section~G encodes mission-shaped situations that random generation is
unlikely to produce: a single bridge edge that any feasible solution must
keep, a mission-critical subgraph whose connectivity dominates the
objective, and stale edges whose freshness penalty must price them out.
Each scenario encodes a domain intuition as a checkable assertion.

\paragraph{H: solver benchmark on the real $Q$ matrix.}
Section~H is the scientific core of the solver audit. Four serious solvers
are run on the actual $Q$ matrices the project produces, across 25
instances with enumerated ground truth: our bit-level simulated annealer
(56\% exact), D-Wave's neal simulated annealer (60\% exact), the tabu
search from the D-Wave toolchain (64\% exact), and simulated quantum
annealing (48\% exact). The exact linearized CBC solve supplies ground
truth and is cross-checked against enumeration on instances up to 60 QUBO
variables. The production projected annealer is deliberately excluded from
this comparison, because it exploits problem structure the other solvers do
not see and would contaminate a like-for-like audit; it appears only as an
engineering note, held to a 90\% exactness floor. This exclusion is itself
part of the methodology: the audit would rather under-claim than compare
unlike objects.

\paragraph{I: trust-GNN layer.}
Section~I checks the message-passing layer in isolation: adjacency
weighting on the pruned graph, boundedness of hidden states under the
sigmoid nonlinearity, correct behavior of the scoring head, and the
capacity-aware greedy assignment rule, including its behavior when
capacities are exhausted.

\paragraph{J: optimality-gap benchmark.}
Section~J quantifies how far the production solver sits from proven optima
using the normalized gap $(H_{\mathrm{SA}} - H^{*}) / \max(1, |H^{*}|)$,
which remains well defined even when energies cross zero. Small instances
are measured against enumeration, where the projected annealer is 100\%
exact, and larger instances, up to 817 binary variables, are measured
against MILP or labeled best-known solutions.

\paragraph{K: simulation harness.}
Section~K audits the auditor's instruments: the causal simulator, the
outcome sampling, the trust updates, and the paired statistical machinery
(bootstrap confidence intervals and permutation tests) are exercised on
constructed cases with known answers, so that a bug in the measurement
apparatus cannot masquerade as a result.

\paragraph{L: review regressions.}
Section~L is a growing graveyard of past wounds. Every legitimate issue
raised in the earlier external review round is encoded as a permanent
regression check, including the requirement that mission success depend
causally on the pruned topology rather than being correlated with it after
the fact.

\paragraph{M: MILP ground truth against the projected annealer.}
Section~M compares the production solver against the exact MILP on twenty
larger instances. The measured outcome, 95\% exact with a mean true gap of
$7\times10^{-5}$, is enforced with an exactness floor of 80\%, so a
degradation of the production solver fails the build even if it remains
competitive in absolute terms.

\paragraph{N: sensitivity diagnostics.}
Section~N perturbs the annealer's temperatures, cooling rates, restart
counts, penalty scales, and pair-move probability, and requires the solver
to remain 87\% exact across the entire perturbation grid. This section
exists to prove that the headline numbers are not the product of a lucky
hyperparameter dial setting: the solver must be good everywhere nearby, not
only at the chosen point.

\subsection{Engine and server tests for the mission console}

The browser-based mission console is a separate software system with its
own failure modes, and it carries its own test layers, run through the
JavaScript test runner for the engine and Python unit tests for the
server. The engine tests are physics-first: they verify the mathematics of
motion before anything visual is trusted.

The differential-drive integrator is verified in closed form. For constant
wheel speeds the true trajectory is a circular arc, and the engine's
integrator is checked against the analytic arc to a tolerance of
$10^{-9}$, including provably singularity-free behavior at zero yaw rate,
where the arc degenerates to a straight line and a naive implementation
divides by zero. Wheel odometry is checked kinematically: over a 200-step
numerical trajectory, each wheel's contact point must travel exactly its
rim distance along the arc, tying the rendered motion to the commanded
wheel speeds. Speed saturation is required to preserve curvature, so that a
saturated rover follows the same arc more slowly rather than cutting the
corner. Spawn layout generation must be deterministic and collision-free
for fleet sizes from 2 to 16 rovers across seeds. A 1800-frame soak test
across four configurations asserts two invariants on every frame: no two
rovers ever overlap, and no rover ever escapes the arena. A 16-rover,
40-second deadlock test requires at least 14 rovers to keep moving,
guarding against a fleet that politely grids itself to a standstill.

The graph and configuration layers are tested with the same severity:
malformed planner data must be rejected without corrupting engine state;
the widest-path tree used for topology display must be correct against
reference computations; parameter updates must snap to their legal grids;
every shipped preset must validate against the parameter specification; and
the bootstrap confidence interval computation must be bit-for-bit
deterministic for a fixed seed.

The server tests cover the HTTP boundary and the session machinery. Request
validation rejects NaN, infinity, out-of-range values, and wrongly typed
fields, and malformed bodies produce clean 400 responses rather than stack
traces. The session store enforces its capacity limit and evicts correctly.
The paired-comparison machinery is checked for its essential fairness
property: follower sessions must share the primary session's kept topology,
fault events must be delivered to all followers, and solver baselines must
run on the identical $Q$ matrix as the primary solve, so that a comparison
panel can never win by secretly solving an easier problem.

\subsection{What this methodology buys}

The practical consequence of this discipline is that every quantitative
claim in the results chapter is backed by at least one check that would
fail if the number regressed. The 56\%, 60\%, 64\%, and 48\% exactness
rates of the external solvers, the 95\% exactness and $7\times10^{-5}$
mean gap of the production annealer, the 87\% robustness across
perturbations, and the statistical margins of the ablation study are not
measurements taken once and transcribed into prose; they are floors and
assertions that the build enforces on every run. Equally important, the
methodology buys the right to report negative results with confidence. When
we state that the QUBO pruner loses to a simpler threshold pruner on raw
mission success at ten-rover scale, or that readout training yields no
significant gain, those statements survive the same hostile scrutiny as the
positive ones. An audit suite that only protected flattering numbers would
be marketing; one that protects every number equally is evidence.
